In the past, Industrial Control Systems (ICS) were kept isolated from external networks. However, the rapid convergence of Information Technology (IT) and Operational Technology (OT) has exposed these environments to modern cyber threats. Because legacy industrial protocols like Modbus TCP lack built-in security features such as authentication or encryption, they fail to meet modern cybersecurity standards, leaving critical infrastructure highly vulnerable. Attackers can easily exploit these protocol weaknesses to inject malicious commands, disrupt physical processes, or bypass boundary defenses, directly compromising operational safety. Since industrial networks demand extreme real-time reliability, deploying an effective monitoring solution without adding latency or processing overhead is a critical challenge in OT security.

To address these issues, this thesis presents a Modbus TCP threat monitoring framework that bridges technical detection with international security standards. The design integrates the ISA/IEC 62443-3-3 Foundational Requirements (FRs) and the ICS Cyber Kill Chain to systematically identify abnormal traffic and adversary behaviors. We built a lightweight, passive monitoring system using the open-source engines Suricata and Graylog. By leveraging deep packet inspection (DPI), custom detection rules, and SIEM correlation, the framework converts raw packet-level anomalies into actionable, context-rich security alerts. Because our solution relies entirely on open-source tools, it drastically reduces deployment and maintenance overhead, making it highly accessible for organizations with tighter OT security budgets.

To streamline incident response, the system triggers real-time alerts via Slack webhooks as soon as a threat is detected. Furthermore, alerts are mapped directly to the MITRE ATT\&CK for ICS framework, giving operators immediate, structured context regarding the attacker’s tactical stage and technical objectives.

This study validates the defensive capabilities of the proposed framework through proof-of-concept (PoC) experiments in a virtualized laboratory platform, covering scenarios such as reconnaissance scanning, protocol fuzzing, unauthorized parameter tampering, critical diagnostic command injection, and denial-of-service (DoS) flooding. The empirical results demonstrate that the system achieves high alert timeliness, maintaining an average end-to-end alert latency of under 10 seconds across all attack scenarios. Under high-throughput loading during peak-rate stress testing, the passive defense pipeline successfully ingested a sustained flooding volume of up to 29,128.00 Packets Per Second (PPS) (approximately 13.36 Mbps), achieving a 0.00\% packet drop rate at the Suricata kernel socket level and a 0.00\% log loss rate in the Graylog SIEM layer. Furthermore, supplementary cross-validation using two third-party public industrial datasets (CICModbusDataset2023 and the 4SICS Geek Lounge traffic) confirms the portability and detection performance of the custom rulesets within heterogeneous external OT network environments.

Overall, our passive framework significantly enhances security visibility, alert timeliness, and threat traceability in industrial control networks without interfering with or delaying real-time Modbus TCP communications. It proves to be a highly viable, production-ready solution that can be reliably deployed on lightweight, consumer-grade hardware.

\noindent
{Keywords:  Modbus TCP、ISA/IEC 62443-3-3、ICS Cyber Kill Chain、Suricata、Graylog、MITRE ATT\&CK for ICS}